% ===================================================================== INTRODUCTION
\section{Introduction}
\label{sec:intro}

Permissioned blockchains may need to replace a consensus engine to change their fault model,
finality semantics, or operating regime. Even with an identical validator committee on both sides,
this is not a routine software upgrade: source and target engines can recognize different quorums
and different points of irreversibility. If validators transfer authority at different boundaries,
one may authorize source history while another commits an incompatible target history. The
process controls in \cref{sec:rq2} exercise threshold-induced target conflicts and gate
authorization, but do not reproduce the full cross-domain source-retirement adversary.

One operational answer is to halt service for a controlled transition. Hyperledger Fabric's
documented Raft-to-BFT migration enters maintenance mode, rejects normal transactions, restarts
ordering nodes, and does not permit a return to Raft after BFT begins committing
\cite{fabric,fabricbftmigration}. Runtime switching offers a different point of comparison:
BFTBrain selects among BFT protocols in one cluster~\cite{bftbrain}, Cox switches among BFT
algorithms with checkpoint and recovery mechanisms~\cite{cox}, and ATBFT alternates BFT
modes~\cite{atbft}. These systems do not reduce to a generic quorum-threshold change, and their
handoff mechanisms deserve comparison on their own terms (\cref{sec:related}). Our question is
more specific: for one fixed, identical committee, how can a source remain authoritative while a
target checks finalized state, and how can the subsequent boundary, source retirement, and
provisional target authority be certified? The handoff deliberately pauses finalization while
certificates are collected; it is not uninterrupted consensus.

\paragraph{Boundary agreement and source retirement}
Consider one committee $\VM$ of $n$ validators, at most $f$ of them Byzantine, with $n\ge 3f+1$. The
first obligation concerns the \emph{boundary certificate}. A certificate of $q$ matching attestations
must be formable when $f$ validators are silent, so $q\le n-f$. Two conflicting certificates must
share a correct signer, so $2q-n>f$. Together these give
\begin{equation}
  \tfrac{n+f}{2} \;<\; q \;\le\; n-f .
  \label{eq:intro-window}
\end{equation}
For an arbitrary-subset boundary certificate, the familiar threshold $2f+1$ lies in this window
only when $n=3f+1$ (\cref{fig:quorum}, panel~a). On a larger committee it no longer guarantees a
correct common signer. This does not classify protocols with restricted voting sets or different
fault assumptions. \sage{} chooses $q=n-f$, the largest available threshold, not the only safe one.

\begin{figure*}[t]
\centering
\includegraphics[width=\textwidth]{figures/fig_quorum.pdf}
\caption{Quorum geometry. (a)~Safe and available boundary thresholds by committee.
(b)~Fence-only exclusion contract at $n=7$, $f=2$. The circled $(4,5)$ point is an inadmissible
source policy, not a counterexample to the locked \sage{} design; the square marks $(5,5)$.}
\label{fig:quorum}
\end{figure*}

A certificate that records agreement but imposes no source-voting restriction does not by itself
transfer authority. To characterize this case separately, let $q_1$ be a source finalization quorum
that could remain active and $q_F$ the quorum that signs a durable \emph{source fence}. The
\emph{fence-only exclusion contract} is
\begin{equation}
  q_F + q_1 \;>\; n + f ,
  \label{eq:intro-fence}
\end{equation}
which forces every fence quorum and every admitted residual source quorum to share a correct
validator. \Cref{th:authority-exclusion} shows that \eqref{eq:intro-fence} excludes residual source
finalization even when the fence certificate reaches only a target quorum.
\Cref{th:fence-admission} shows that such a fence is available if and only if $q_1>2f$, with $q_F=n-f$
always sufficient. The rule is sharp. At $(n,f,q_1,q_F)=(7,2,4,5)$ equality holds, so a majority-based
source is rejected at admission rather than silently deployed (\cref{fig:quorum}, panel~b).
The two inequalities are logically independent over arbitrary parameter tuples
(\cref{prop:independence}); this does not establish that two separate certificates are necessary
inside every migration protocol.

In the \emph{locked} \sage{} design, each readiness signer durably stops source authorization
above its candidate boundary before releasing its share. An $n-f$ \CutCert{} then retires at least
$n-2f$ correct signers, leaving at most $2f<q_1$ possible source voters. It already guarantees
source quiescence, including after restart (\cref{lem:fence-live}). \FenceCert{} is retained as
explicit evidence tying durable retirement to the complete manifest, engine generation, and admitted
policies, not as an additional safety necessity for that locked profile
(\cref{prop:retirement-evidence}). This dependency is part of the design rather than a claim that
removing each certificate must cause a fork.

\paragraph{Design}
\sage{} permits at most one finalizing engine throughout the handoff. During dual-run, the legacy engine
$\Csrc$ remains the only finalizer. The target $\Ctgt$ re-executes the same finalized blocks without
proposing, so divergence becomes observable before target decisions can bind the chain. After a
stability window, validators attest to one boundary tuple, and $n-f$ matching attestations form a
\CutCert. Separate certificate domains then take over the remaining obligations. An $n-f$ \ManCert{}
authenticates the complete migration manifest. A $q_F$ \FenceCert{} provides policy-bound evidence
that at least $q_F-f$ correct signers durably installed the specified source fence
(\cref{prop:retirement-evidence}); it does not imply installation by every validator.
Target authority requires both certificates and a local durable fence. \AbortCert{} and \SealCert{} use a
terminal authority domain. An abort before the deadline $\hr$ restores the certified anchor under a
\emph{fresh} source generation and never revives the fenced one. Only a seal promotes the target
prefix to absolute finality.

\paragraph{Contributions}
This paper makes four contributions.
\begin{enumerate}[leftmargin=*,label=\textbf{C\arabic*.}]
  \item \textbf{Migration-specific obligations, not new quorum arithmetic.} For a fixed committee,
  we distinguish agreement on an arbitrary-subset boundary \eqref{eq:intro-window} from exclusion
  of residual source authority under a fence-only contract \eqref{eq:intro-fence}. Both inequalities
  instantiate established Byzantine-quorum intersection~\cite{byzantinequorums}. The contribution
  is to identify which obligation each migration state and certificate witnesses: under durable
  Rule~\ref{rule:lock} and the admitted $q_1>2f$ policy, \CutCert{} signers already stop source
  progress, whereas \FenceCert{} certifies later manifest- and policy-bound retirement by its
  correct signers. We do not claim that the latter is independently required for source exclusion
  in this locked design.
  \item \textbf{A certified cross-engine handoff design.} Under an abstract engine interface,
  \sage{} retains one finalizer during shadow execution, commits a boundary \CutCert{},
  authenticates the complete manifest, records durable source fencing, and admits provisional
  target finality only after the specified certificate and local-fence gates. Certified abort
  returns to the anchor in a fresh source generation; certified seal makes the target prefix
  absolute. Signed receipts carry the client-visible finality tier.
  \item \textbf{Bounded formal evidence with falsifying controls.} TLA+ models checked with TLC
  exercise partial certificate delivery, per-validator knowledge, durable fence installation,
  crash--restart, and Byzantine participation in both engines. Separate finite models examine
  fixed-candidate agreement at $n\in\{4,7,10\}$, manifest agreement, and absolute-finality
  reversion. Broken controls are falsified by TLC; these finite checks are not an implementation
  refinement argument.
  \item \textbf{Separated cost and authorization evidence.} A deterministic simulator reports
  event-time costs and parameter sensitivity. Concurrent TCP processes in a PoA-to-HotStuff
  subset exercise constructed threshold conflicts, gate refusal, crash--restart, and convergence
  after a partition. They do not execute the complete networked certificate and terminal path,
  or an authentic external switching baseline.
\end{enumerate}

Here ``live migration'' means continued source service during preparation followed by a certification
pause, not uninterrupted block production. The pause bound is conditional on common readiness,
certificate delivery, durable-write bounds, and target liveness (\cref{th:liveness}). The evaluated
network path is a \CutCert-gated local-fence subset of the design, with a crash-only source;
committee-wide manifest, fence, and terminal stages are not evaluated end to end. Neither its
simulator timings nor its constructed process controls establish full-protocol performance.

The rest of the paper is organized as follows. \Cref{sec:related} positions \sage{} against
reconfiguration, deployed transitions, and runtime switching. \Cref{sec:model} fixes the system and
threat model, \cref{sec:protocol} specifies the protocol, and \cref{sec:analysis} proves its
properties. \Cref{sec:eval} reports the evaluation, \cref{sec:discussion} discusses generality and
scope, and \cref{sec:conclusion} concludes.
